
\documentclass[11pt]{article}

\usepackage{amsmath}
\usepackage{amssymb}
\usepackage{graphicx}
\usepackage{hyperref}
\usepackage{geometry}
\usepackage{booktabs}
\geometry{margin=1in}

\title{The Prime Exponent Tree (PET): \\
A Structural View of Integer Factorization}

\author{Giancarlo Comneno \\
\small Independent Research}

\date{\today}

\begin{document}

\maketitle

\begin{abstract}
We introduce the \textbf{Prime Exponent Tree (PET)}, a recursive tree representation
of integers derived from prime factorization. In this model, every integer is
represented as a tree whose nodes correspond to prime bases and whose children
recursively encode the structure of the corresponding exponents.

Using a complete dataset of all integers up to $10^6$, we empirically explore the
space of PET structures. Surprisingly, despite the large number of integers,
only a small number of distinct structural shapes appear. Among the first one
million integers, we observe only 78 distinct structural shapes, and the
distribution of these shapes is highly concentrated: seven shapes account for
approximately 87\% of all integers.

We further measure the structural entropy of the PET distribution and observe
that the complexity of integer multiplicative structure grows approximately as
$\log\log N$.
\end{abstract}

\section{Introduction}

Integers grow rapidly, but their multiplicative structure evolves much more
slowly. Classical results such as the Hardy--Ramanujan theorem show that the
number of distinct prime factors of an integer grows on average like
$\log\log n$.

The goal of this work is to study the \emph{structure} of integer
factorizations using a recursive tree representation called the
Prime Exponent Tree (PET).

\section{The Prime Exponent Tree}

Let $N \ge 2$ with prime factorization

\begin{equation}
N = \prod_i p_i^{e_i}.
\end{equation}

The PET of $N$ is defined recursively:

\begin{itemize}
\item each prime $p_i$ becomes a node
\item if $e_i = 1$ the node is a leaf
\item if $e_i > 1$ the exponent is represented by a PET subtree
\end{itemize}

Example:

\begin{verbatim}
12 = 2^2 * 3

•
├─2
│  •
│  └─2
└─3
\end{verbatim}

\section{Dataset}

We generated a dataset containing the PET representation of all integers
from $2$ to $1{,}000{,}000$.

Each record contains:

\begin{itemize}
\item the integer $n$
\item the PET structure
\item derived structural metrics
\end{itemize}

The dataset contains 999,999 integers.

\section{Structural Shape Distribution}

We define the \emph{shape} of a PET tree as the structure obtained by
removing the specific prime labels while preserving the recursive topology.

Surprisingly, among the first one million integers we observe only:

\begin{center}
\textbf{78 distinct structural shapes}.
\end{center}

This indicates an extremely strong compression of the multiplicative
structure of integers.

\subsection{Dominant Shapes}

The distribution of shapes is highly concentrated.
The seven most common structures account for approximately 87\% of integers.

\begin{center}
\begin{tabular}{ccc}
\toprule
Structure & Interpretation & Approx Frequency \\
\midrule
$pq$ & two primes & $\sim21\%$ \\
$pqr$ & three primes & $\sim21\%$ \\
$p^2qr$ & square times two primes & $\sim13\%$ \\
$pqrs$ & four primes & $\sim9\%$ \\
$p^2qrs$ & square plus three primes & $\sim8\%$ \\
$p$ & prime & $\sim8\%$ \\
$p^2q$ & square times prime & $\sim7\%$ \\
\bottomrule
\end{tabular}
\end{center}

\section{Height Distribution}

The PET height measures recursive depth of exponent structure.

Empirical distribution for integers $\le 10^6$:

\begin{center}
\begin{tabular}{ccc}
\toprule
Height & Count & Percentage \\
\midrule
1 & 607925 & 60.79\% \\
2 & 347997 & 34.80\% \\
3 & 44069 & 4.41\% \\
4 & 8 & $\approx 0\%$ \\
\bottomrule
\end{tabular}
\end{center}

Thus:

\begin{center}
\textbf{95.6\% of integers have PET height $\le 2$}.
\end{center}

\section{Branching Distribution}

The number of children of the root corresponds to $\omega(n)$,
the number of distinct prime factors.

\begin{center}
\begin{tabular}{ccc}
\toprule
$\omega(n)$ & Count & Percentage \\
\midrule
1 & 78734 & 7.87\% \\
2 & 288726 & 28.87\% \\
3 & 379720 & 37.97\% \\
4 & 208034 & 20.80\% \\
5 & 42492 & 4.25\% \\
6 & 2285 & 0.23\% \\
7 & 8 & $\approx 0\%$ \\
\bottomrule
\end{tabular}
\end{center}

The distribution peaks near $\omega(n)=3$, consistent with
the classical Hardy--Ramanujan theorem.

\section{Structural Entropy}

Let $p_i$ be the probability of shape $i$.
We define structural entropy

\begin{equation}
H = -\sum_i p_i \log p_i.
\end{equation}

For integers up to $10^6$ we obtain

\begin{center}
$H \approx 2.35$
\end{center}

while the maximum possible entropy for 78 shapes is

\begin{center}
$\log 78 \approx 4.36$.
\end{center}

Thus the effective number of structural configurations is

\begin{equation}
e^H \approx 10.5.
\end{equation}

\section{Entropy Growth}

Empirically we observe:

\begin{center}
\begin{tabular}{ccc}
\toprule
$N$ & $H(N)$ & $\log\log N$ \\
\midrule
$10^4$ & 2.18 & 2.22 \\
$10^5$ & 2.28 & 2.44 \\
$10^6$ & 2.35 & 2.63 \\
\bottomrule
\end{tabular}
\end{center}

This suggests the empirical relation

\begin{equation}
H(N) \sim \log\log N.
\end{equation}

\section{Discussion}

Despite the enormous number of integers, their multiplicative
structure appears to occupy a very small combinatorial space.

Most integers fall into a handful of structural categories,
and the complexity of this space grows extremely slowly.

\section{Conclusion}

The Prime Exponent Tree representation reveals a surprisingly
compact structure underlying integer factorizations.

Empirical evidence suggests that the complexity of integer
structure grows roughly as $\log\log N$, indicating that
multiplicative structure is dramatically simpler than
numerical magnitude.

\end{document}
